\documentclass[11pt,a4paper]{article}
\usepackage[utf8]{inputenc}
\usepackage[T1]{fontenc}
\usepackage{amsmath,amssymb}
\usepackage{geometry}
\geometry{margin=2.5cm}
\usepackage{xcolor}
\usepackage[most]{tcolorbox}
\usepackage{fancyhdr}
\usepackage{hyperref}
\hypersetup{colorlinks=true,linkcolor=blue!50!black,urlcolor=blue!50!black,citecolor=blue!50!black}
\usepackage{enumitem}
\usepackage{booktabs}

\definecolor{defcol}{HTML}{2A6F8F}
\definecolor{thmcol}{HTML}{6B4FA0}
\definecolor{excol}{HTML}{2E7D4F}
\definecolor{warncol}{HTML}{C24A1E}
\definecolor{boxbg}{HTML}{F4F6F8}

\newtcolorbox{definition}[1]{colback=boxbg,colframe=defcol,fonttitle=\bfseries,title={Définition~---~#1}}
\newtcolorbox{theoreme}[1]{colback=boxbg,colframe=thmcol,fonttitle=\bfseries,title={Résultat~---~#1}}
\newtcolorbox{postulat}{colback=white,colframe=excol,fonttitle=\bfseries,title=Postulat}
\newtcolorbox{limite}{colback=white,colframe=warncol,fonttitle=\bfseries,title=Limite assumée}
\newtcolorbox{remarque}{colback=boxbg,colframe=black!40,fonttitle=\bfseries\itshape,title=Remarque}

\newcommand{\demo}{\par\smallskip\noindent\textbf{Démonstration / dérivation.}\ }
\newcommand{\qedhere}{\hfill$\blacksquare$\par\smallskip}

\pagestyle{fancy}
\fancyhf{}
\fancyhead[L]{\small Signaux précurseurs performatifs}
\fancyhead[R]{\small\thepage}
\renewcommand{\headrulewidth}{0.3pt}

\title{\textbf{Signaux précurseurs performatifs}\\[4pt]
\large Un premier modèle formel de la réflexivité dans la détection de bascules collectives}
\author{}
\date{Septembre 2026}

\begin{document}
\maketitle

\begin{center}
\begin{minipage}{0.85\textwidth}
\small\itshape
Ce document pose et résout un premier modèle mathématique reliant deux
champs établis séparément — les signaux précurseurs de bascule
(ralentissement critique) et la prédiction performative — sans qu'aucun
travail trouvé ne les combine formellement à ce jour (vérifié par
recherche approfondie puis par recoupement direct des bibliographies des
synthèses les plus récentes de chaque champ). Ce n'est pas présenté
comme une théorie complète de la réflexivité, mais comme un premier
modèle jouet, dérivé et vérifié numériquement, destiné à servir de point
de départ à une collaboration ouverte plutôt que comme un résultat
définitif.
\end{minipage}
\end{center}

\tableofcontents

\section{Motivation et positionnement honnête}

Les indicateurs de ralentissement critique (variance, autocorrélation)
supposent un paramètre de contrôle $\lambda(t)$ purement exogène — le
système physique ou écologique observé ne «~lit~» jamais
l'indicateur calculé sur lui-même. Cette hypothèse est raisonnable pour
un lac ou un plasma. Elle est fausse par construction pour un système
social observé par un outil comme Hélios : si l'alerte est publiée, les
acteurs peuvent y réagir, et cette réaction change la trajectoire même
que l'indicateur cherchait à décrire.

\begin{remarque}
Ce problème porte un nom en apprentissage automatique — la
\emph{prédiction performative} (Perdomo, Zrnic, Mendler-Dünner \& Hardt,
2020) — et un nom informel en économie/finance — la \emph{réflexivité}
(Soros) ou la \emph{critique de Lucas}. Mais aucun des deux cadres n'a
été formalisé dans le langage mathématique précis des signaux
précurseurs (variance, autocorrélation, ralentissement critique). C'est
exactly le vide que ce document tente de combler, modestement.
\end{remarque}

\section{État de l'art et positionnement précis}

Une revue approfondie (recherche multi-sources, puis vérification
directe des bibliographies des deux synthèses les plus récentes de
chaque champ — Hardt \& Mendler-Dünner, 2025, \emph{Statistical
Science} ; «~Dissecting Performative Prediction~», 2026) ne trouve
aucun travail combinant formellement les deux cadres. Les travaux les
plus proches, chacun d'un seul côté du pont :

\begin{itemize}[leftmargin=1.4em]
\item \textbf{Diekert, Heyen, Nesje \& Shayegh (2025)}, \emph{J. R. Soc.
  Interface} — le lien formel le plus proche : un système d'alerte
  précoce change les décisions d'un planificateur, et peut produire un
  comportement \emph{plus} risqué en l'absence d'alerte. Un seul
  décideur rationnel, pas de population réactive, pas de statistique de
  ralentissement critique.
\item \textbf{Bauch, Sigdel, Pharaon \& Anand (2016)}, \emph{PNAS} — la
  rétroaction humaine peut «~atténuer le signal précurseur~»,
  rendant une bascule à venir plus difficile à détecter, dans un modèle
  couplé forêt-opinion.
\item \textbf{Diks, Hommes \& Wang (2019)}, \emph{Empirical Economics} —
  notent, en une phrase, qu'une crise prédite «~peut causer une
  réaction de panique des agents et ainsi déclencher une crise à
  venir encore plus tôt~», sans le formaliser.
\item \textbf{Sornette (2003)}, \emph{Physics Reports} — liste, de
  façon informelle, les scénarios où une prédiction publique de krach
  s'auto-réalise ou s'auto-invalide.
\item \textbf{Drehmann \& Juselius (2014)}, BIS/\emph{Int. J.
  Forecasting} — notent qu'un indicateur d'alerte utilisé en politique
  publique «~deviendra sujet à la critique de Lucas usuelle~», et
  que la perte de pouvoir prédictif serait «~plutôt un signe de son
  succès~».
\end{itemize}

\begin{limite}
Cette revue s'appuie sur une recherche documentaire, pas sur une base de
données de citations exhaustive. L'absence de résultat n'est pas une
preuve d'absence définitive.
\end{limite}

\section{Le modèle}

\subsection{Rappel du cadre sans rétroaction}

Le modèle de ralentissement critique déjà établi (processus
d'Ornstein-Uhlenbeck près d'une bifurcation) s'écrit :
$$dy = \lambda(t)\,y\,dt + \sigma\,dW, \qquad \lambda(t) \to 0^-$$
avec la variance stationnaire $\text{Var}(y) = \sigma^2/(2|\lambda|)$ et
l'autocorrélation $\text{AC1} \to 1$ quand $\lambda\to0^-$ — le signal
précurseur classique. Ce résultat suppose $\lambda(t)$ exogène.

\subsection{Introduire la performativité}

\begin{postulat}
Un détecteur observe la variance glissante $V_t$ et publie une alerte
$a_t = A(V_t) \in [0,1]$, une fonction sigmoïde croissante de $V_t$.
Cette alerte modifie en retour le paramètre de stabilité effectif du
système :
$$\lambda_{\text{eff}}(t) = \lambda(t) + \kappa \cdot A(V_t), \qquad A(V) = \frac{1}{1+e^{-k(V-V_c)}}$$
où $\kappa$ est le \emph{coefficient de rétroaction performative}
($\kappa<0$ : réponse stabilisatrice/auto-invalidante ; $\kappa>0$ :
réponse déstabilisatrice/auto-réalisatrice), $k$ la sensibilité de
l'alerte, $V_c$ le seuil de déclenchement.
\end{postulat}

\begin{remarque}
$\kappa$ joue exactement le rôle de la sensibilité $\varepsilon$ dans le
cadre de Perdomo et al. (le décalage de distribution induit par le
paramètre du modèle), transposé ici au paramètre de stabilité d'un
processus stochastique plutôt qu'à une distribution de caractéristiques.
\end{remarque}

\subsection{L'équation de point fixe}

En régime quasi-stationnaire ($V_t \approx \sigma^2/(2|\lambda_{\text{eff}}|)$,
séparation d'échelles de temps admise), on obtient une équation implicite
auto-cohérente pour $\lambda_{\text{eff}}$ :
$$\lambda_{\text{eff}} = \lambda(t) + \kappa \cdot A\!\left(\frac{\sigma^2}{2|\lambda_{\text{eff}}|}\right)$$

\begin{remarque}
Structure identique à l'équation d'auto-cohérence de Kuramoto déjà
résolue pour H4 (§5.8 du cahier des charges Hélios) — le champ moyen
$r$ y jouait le rôle que joue ici $\lambda_{\text{eff}}$.
\end{remarque}

\section{Analyse du point fixe}

\begin{theoreme}{Deux régimes qualitativement distincts}
Selon le signe de $\kappa$, l'équation de point fixe se comporte de deux
façons radicalement différentes lorsque $\lambda(t)\to0^-$ :
\begin{itemize}[leftmargin=1.4em]
\item $\kappa=0$ : $\lambda_{\text{eff}}=\lambda(t)$, on retrouve le cas
  classique (vérification de cohérence).
\item $\kappa<0$ (auto-invalidant) : $\lambda_{\text{eff}}$ reste
  \textbf{borné loin de zéro} même quand $\lambda(t)\to0^-$ — la
  variance ne diverge plus, elle plafonne.
\item $\kappa>0$ (auto-réalisateur) : il existe un $\lambda_c(\kappa)<0$
  tel que \textbf{le point fixe cesse d'exister} pour
  $\lambda(t)>\lambda_c(\kappa)$ — la bascule (au sens de cette
  équation) survient avant que le système «~nu~» n'y soit arrivé.
\end{itemize}
\end{theoreme}

\demo Résolution numérique par bissection de
$f(\lambda_{\text{eff}}) = \lambda_{\text{eff}} - \lambda(t) - \kappa A(\sigma^2/2|\lambda_{\text{eff}}|) = 0$,
pour $\sigma=1$, $k=10$, $V_c=2$ :

\textbf{Cas $\kappa<0$} — extrait ($\kappa=-0{,}1$) :
\begin{center}
\begin{tabular}{@{}rrr@{}}
\toprule
$\lambda(t)$ & $\lambda_{\text{eff}}$ & $\text{Var}$ \\
\midrule
$-1{,}000$ & $-1{,}000$ & $0{,}50$ \\
$-0{,}100$ & $-0{,}199$ & $2{,}51$ \\
$-0{,}010$ & $-0{,}110$ & $4{,}55$ \\
$-0{,}001$ & $-0{,}101$ & $4{,}95$ \\
\bottomrule
\end{tabular}
\end{center}
Même quand $\lambda(t)$ s'approche de $0$ à $10^{-3}$ près,
$\lambda_{\text{eff}}$ reste proche de $-0{,}1$ — un \textbf{plancher de
stabilité performatif}. La divergence classique de la variance
n'apparaît jamais.

\textbf{Cas $\kappa>0$} — extrait ($\kappa=0{,}05$) :
\begin{center}
\begin{tabular}{@{}rl@{}}
\toprule
$\lambda(t)$ & Solution \\
\midrule
$-0{,}500$ & $\lambda_{\text{eff}}=-0{,}500$ \\
$-0{,}100$ & $\lambda_{\text{eff}}=-0{,}050$ \\
$-0{,}050$ & \textbf{pas de solution} \\
$-0{,}020$ & \textbf{pas de solution} \\
\bottomrule
\end{tabular}
\end{center}
Le point fixe disparaît entre $\lambda=-0{,}1$ et $\lambda=-0{,}05$ —
bien avant que le système nu n'atteigne $\lambda=0$. Plus $\kappa$
augmente, plus $\lambda_c(\kappa)$ s'éloigne de $0$ (bascule anticipée
plus tôt). \qedhere

\section{Validation par simulation stochastique complète}

L'analyse précédente repose sur l'approximation quasi-stationnaire. Pour
vérifier qu'elle tient dans un vrai système bruité, avec une variance
glissante calculée en temps réel sur une fenêtre finie (pas la formule
fermée), le système complet a été simulé par Euler-Maruyama :
$$y_{n+1} = y_n + dt\,(\lambda(t_n) + \kappa A(V_n))\,y_n + \sigma\sqrt{dt}\,\xi_n$$

\begin{remarque}[Erreur trouvée et corrigée en cours de route]
Un premier essai avec une fenêtre glissante de $0{,}5$ unité de temps
n'a montré \textbf{aucune différence entre les valeurs de $\kappa$} — le
temps de relaxation du système près de la bifurcation ($\sim 1/|\lambda|$)
dépasse largement la taille de la fenêtre, donc la variance mesurée
restait toujours proche de $0$ et l'alerte ne se déclenchait jamais,
quel que soit $\kappa$. Corrigé en portant la fenêtre à $20$ unités de
temps (recalibrée en vérifiant que la variance mesurée suit bien la
prédiction théorique en régime $\kappa=0$ avant de comparer les régimes
entre eux) — exactement le type d'erreur que ce projet a appris à
détecter systématiquement plutôt qu'à laisser passer.
\end{remarque}

\textbf{Résultat} (\textbf{500 réalisations indépendantes} par valeur de
$\kappa$, graine de base explicite $=42$ — reproductible, pas de
dérivation informelle par fonction de hachage —, $\lambda_0=-1$, taux de
dérive $0{,}002$, seuil de bascule $|y|>8$) :

\begin{center}
\begin{tabular}{@{}lrrr@{}}
\toprule
Régime & Bascules & Taux & Temps moyen (écart-type) \\
\midrule
$\kappa=0$ (classique) & 317/500 & 63,4\% & 504,2 (15,7) \\
$\kappa=-0{,}05$ & 162/500 & 32,4\% & 508,5 (14,3) \\
$\kappa=-0{,}10$ & 51/500 & 10,2\% & 510,9 (13,6) \\
$\kappa=-0{,}20$ & 15/500 & \textbf{3,0\%} & 509,6 (12,6) \\
$\kappa=+0{,}02$ & 373/500 & 74,6\% & 499,3 (16,9) \\
$\kappa=+0{,}05$ & 442/500 & 88,4\% & 493,0 (19,1) \\
$\kappa=+0{,}08$ & 480/500 & 96,0\% & 483,9 (20,8) \\
\bottomrule
\end{tabular}
\end{center}

\begin{remarque}[Vérification de l'absence de bistabilité cachée]
L'analyse du \S4 repose sur une résolution par bissection, qui suppose
implicitement une seule racine dans l'intervalle testé. Pour écarter
l'hypothèse d'un pli à deux solutions (bistabilité/hystérésis) que la
bissection aurait pu manquer, la fonction résiduelle a été scannée
finement (jusqu'à 20 000 points, avec un zoom à 21 points entre
$\lambda=-0{,}06$ et $-0{,}05$) autour de la zone de transition pour
$\kappa=0{,}05$. Résultat : exactement une racine jusqu'à
$\lambda\approx-0{,}0505$, puis zéro — un vrai pli net (tangence),
aucune fenêtre à deux solutions cachée. La caractérisation du \S4 est
donc confirmée, pas seulement supposée.
\end{remarque}

\subsection{Sensibilité aux paramètres de la fonction d'alerte}

Le résultat qualitatif (monotonie du taux de bascule en fonction de
$\kappa$) a été revérifié à \textbf{1000 réalisations}, pour 5
combinaisons de la pente $k$ et du seuil $V_c$ de la fonction d'alerte
($k\in\{5,10,20\}$, $V_c\in\{1{,}5;\,2;\,2{,}5\}$) : la monotonie tient
dans les 5 cas.

\begin{remarque}[Nuance nécessaire — monotonie au sens large, pas stricte partout]
Un test de sensibilité complémentaire au \textbf{seuil de bascule
lui-même} ($|y|>y_{\text{thresh}}$, jusqu'ici jamais testé) montre que la
monotonie stricte ne tient qu'à $y_{\text{thresh}}=8$. À
$y_{\text{thresh}}=12$ (seuil élevé), $\kappa=-0{,}1$ et $\kappa=-0{,}2$
donnent tous deux $0{,}0\%$ — un effet de plancher (les deux régimes
empêchent totalement la bascule dans l'horizon simulé), pas une
inversion de tendance, mais qui rompt la stricte monotonie affichée sans
nuance dans une première version de ce document. Le résultat correct à
retenir : \textbf{monotonie au sens large (jamais d'inversion), avec
saturation possible aux seuils extrêmes} — pas une monotonie stricte
inconditionnelle.
\begin{center}
\begin{tabular}{@{}lrrrrr@{}}
\toprule
$y_{\text{thresh}}$ & $\kappa{=}0$ & $\kappa{=}{-}0{,}1$ & $\kappa{=}{-}0{,}2$ & $\kappa{=}{+}0{,}05$ & $\kappa{=}{+}0{,}08$ \\
\midrule
5  & 97,8\% & 83,4\% & 72,8\% & 99,6\% & 99,6\% \\
8  & 63,4\% & 10,2\% & 3,0\%  & 88,4\% & 96,0\% \\
12 & 35,8\% & 0,0\%  & 0,0\%  & 82,4\% & 93,8\% \\
\bottomrule
\end{tabular}
\end{center}
\end{remarque}

\subsection{Extension au modèle de Kuramoto (H4)}

Pour vérifier que le mécanisme n'est pas propre au modèle nœud-col, la
même construction (alerte $A(r_t)$ sur le paramètre d'ordre $r$ au lieu
de la variance, $K_{\text{eff}}(t) = K_0(t) + \kappa A(r_t)$) a été
appliquée au modèle de synchronisation déjà utilisé pour H4 (§5.8) :

\begin{center}
\begin{tabular}{@{}lrr@{}}
\toprule
Régime & Bascules (sur 500) & Temps moyen \\
\midrule
$\kappa=0$ (classique) & 500/500 & 141,7 \\
$\kappa=-0{,}5$ & 500/500 & 273,0 \\
$\kappa=-1{,}0$ & 495/500 & 402,1 \\
$\kappa=-2{,}0$ & \textbf{0/500} & — \\
$\kappa=+0{,}5$ & 500/500 & 50,5 \\
$\kappa=+1{,}0$ & 500/500 & 22,5 \\
$\kappa=+2{,}0$ & 500/500 & 12,2 \\
\bottomrule
\end{tabular}
\end{center}

\begin{remarque}[Asymétrie trouvée lors d'une relecture critique — à ne pas passer sous silence]
Contrairement à ce qu'une première lecture du tableau ci-dessus suggère,
les deux régimes ne sont \textbf{pas confirmés symétriquement} sur ce
modèle. Un test de sensibilité au seuil de bascule
$r_{\text{tip}}\in\{0{,}6;\,0{,}7;\,0{,}8\}$ (200 réalisations, régime
auto-réalisateur) montre que le \emph{taux} de bascule sature à $100\%$
dès $\kappa=+1$, quel que soit $r_{\text{tip}}$ testé — seul le
\emph{délai moyen} continue à discriminer entre les valeurs de $\kappa$
côté auto-réalisateur. Le mécanisme auto-invalidant, lui, discrimine
bien à la fois sur le taux et sur le délai (voir $\kappa=-2$, effondrement
net du taux). \textbf{À retenir} : sur Kuramoto, c'est la variable
temporelle qui porte l'essentiel du signal auto-réalisateur, pas le taux
de bascule — une asymétrie réelle entre les deux modèles testés, qui
n'existe pas sur le nœud-col (où taux \emph{et} délai discriminent des
deux côtés).
\end{remarque}

Même signature qualitative sur une dynamique de nature complètement
différente (synchronisation collective plutôt qu'une bifurcation
simple) : plancher de stabilité net côté auto-invalidant, accélération
monotone côté auto-réalisateur. Le mécanisme ne dépend donc pas du choix
particulier du modèle sous-jacent.

\begin{limite}
«~0/30 bascule~» à $\kappa=-0{,}2$ signifie \emph{dans l'horizon
temporel simulé} (jusqu'à $\lambda(t)=0{,}05$) — pas que la bascule est
impossible dans l'absolu. Un temps de simulation plus long ou un taux de
dérive de $\lambda(t)$ plus rapide pourrait finir par forcer la bascule
même à $\kappa$ très négatif.
\end{limite}

\section{Limites assumées de ce premier modèle}

\begin{limite}
\begin{itemize}[leftmargin=1.4em]
\item Un seul mécanisme de rétroaction est modélisé (réponse linéaire du
  paramètre de stabilité à une alerte sigmoïde). Les mécanismes de
  \emph{masquage stratégique} (loi de Goodhart — des acteurs qui
  suppriment délibérément les fluctuations mesurées) et d'\emph{aléa
  moral} (comportement plus risqué en l'\emph{absence} d'alerte,
  Diekert et al.) ne sont pas modélisés ici.
\item Le détecteur ne se recalibre pas — pas de boucle de ré-entraînement
  au sens de la «~stabilité performative~» de Perdomo et al.
  (convergence d'une politique de décision sous minimisation de risque
  répétée). Ce modèle teste un $\kappa$ fixé, pas la convergence d'un
  processus d'apprentissage.
\item Validé sur deux modèles déjà internes à Hélios (nœud-col,
  Kuramoto), à 500-1000 réalisations, avec vérification de sensibilité
  aux paramètres de la fonction d'alerte — mais toujours aucun test sur
  données réelles, aucune relecture par un chercheur extérieur au
  domaine, et une vérification bibliographique par recoupement de
  synthèses plutôt que par une base de citations exhaustive
  (Semantic Scholar / Google Scholar complet).
\item C'est un premier modèle jouet, pas une théorie de la réflexivité
  sociale. Présenté comme point de départ pour une collaboration plus
  large, pas comme un résultat clos.
\end{itemize}
\end{limite}

\section{Prochaines étapes possibles}

\begin{enumerate}[leftmargin=1.4em]
\item Formaliser le mécanisme de masquage stratégique (Goodhart) comme
  second modèle, complémentaire à celui-ci.
\item Étudier la convergence d'un détecteur qui recalibre son seuil
  $V_c$ à partir des données déjà influencées par ses propres alertes
  passées — la vraie notion de «~stabilité performative~» au sens
  de Perdomo et al., pas seulement un $\kappa$ fixe.
\item Tester sur un vrai jeu de données où la réflexivité est plausible
  et où des signaux de type ralentissement critique ont déjà été
  cherchés — les marchés financiers (Diks, Hommes \& Wang, 2019 ;
  Guttal et al., 2016) sont le candidat le plus direct.
\item Décision ouverte : héberger ce travail comme extension d'Hélios,
  ou comme projet et interface séparés, étant donné sa nature plus
  spéculative et son statut de première tentative sur un vide de
  recherche identifié plutôt que d'application d'une méthode établie.
\end{enumerate}

\section*{Références}
\begin{itemize}[leftmargin=1.2em]
\item Perdomo, J., Zrnic, T., Mendler-Dünner, C., \& Hardt, M. (2020). "Performative Prediction." *ICML 2020*, PMLR 119.
\item Hardt, M., \& Mendler-Dünner, C. (2025). "Performative Prediction: Past and Future." *Statistical Science*, 40(3), 417–436.
\item Diekert, F., Heyen, D., Nesje, F., \& Shayegh, S. (2025). "Do early warning signals of tipping points lead to better decisions?" *Journal of the Royal Society Interface*, 22(225), 20240864.
\item Bauch, C. T., Sigdel, R., Pharaon, J., \& Anand, M. (2016). "Early warning signals of regime shifts in coupled human–environment systems." *PNAS*, 113(51), 14560–14567.
\item Diks, C., Hommes, C., \& Wang, J. (2019). "Critical slowing down as an early warning signal for financial crises?" *Empirical Economics*, 57(4), 1201–1228.
\item Sornette, D. (2003). "Critical market crashes." *Physics Reports*, 378, 1–98.
\item Drehmann, M., \& Juselius, M. (2014). "Evaluating early warning indicators of banking crises: Satisfying policy requirements." *International Journal of Forecasting*, 30(3).
\item Scheffer, M. et al. (2009). "Early-warning signals for critical transitions." *Nature*, 461, 53–59.
\item Guttal, V., Raghavendra, S., Goel, N., \& Hoarau, Q. (2016). "Lack of Critical Slowing Down Suggests that Financial Meltdowns Are Not Critical Transitions, yet Rising Variability Could Signal Systemic Risk." *PLoS ONE*, 11(1), e0144198.
\end{itemize}

\end{document}
